\section{Contamination Channel Capacity}
\label{sec:capacity}

\subsection{Setup}

Consider a sequence of $n$ processing operations
$f_1, f_2, \dotsc, f_n \in \mathrm{Hom}(\SC)$. Each operation acts on
the (random) presence of a haram-bit $X \in \{0,1\}$ as a discrete
memoryless channel $W_i \colon X \mapsto Y_i$ with transition matrix
$\Pi_i$. Empirical estimates of $\Pi_i$ for canonical operations
(centrifugation, distillation, washing cycles, etc.) are tabulated in
\Cref{sec:eval}.

\subsection{Capacity}

\begin{definition}[Contamination channel capacity]
\label{def:capacity}
For a chain $W = W_n \circ \dotsb \circ W_1$ define
\[
   \Ch(W) \;=\; \max_{p_X} I(X; Y_n)
\]
where the maximum ranges over input distributions on
$\{0,1\}$ subject to the operational admissibility constraints inherited
from $\SC$.
\end{definition}

\begin{theorem}[Survival bound]
\label{thm:survival}
Under \cref{def:capacity},
\[
   \Pr[\, Y_n = 1 \mid X = 1 \,]
   \;\leq\; 2^{\,\Ch(W) - n \cdot \delta(W)}
\]
where $\delta(W)$ is the per-step Kullback--Leibler dissipation of $W$.
In particular, $n \geq \Ch(W) / \delta(W) - \log_2 \varepsilon^{-1}$
operations suffice to drive haram-signal survival below any
$\varepsilon > 0$.
\end{theorem}

\Cref{thm:survival} converts the qualitative ``thoroughness'' of a
cleaning protocol into an auditable bit count. The Lean formalization is
in \path{proofs/Ghc/Info.lean}; tabulated $\Pi_i$ for industrial
operations are produced by \path{notebooks/04-channel-estimation.ipynb}.
